\section{总结与延伸}

\subsection{核心技术脉络}

本讲从条件生成的基本问题出发，系统介绍了将扩散模型推向实际应用的关键技术栈：

\begin{table}[H]
\centering
\begin{tabular}{lp{4.5cm}p{4.5cm}}
\toprule
\textbf{技术} & \textbf{解决的问题} & \textbf{核心思想} \\
\midrule
Classifier Guidance & 条件生成 & score = 无条件 score + 分类器梯度 \\
CFG & 去掉外部分类器 & 条件 dropout + 推理时外推 \\
DDIM Inversion & 图像编辑 & 确定性反演 + 换条件重生成 \\
Latent Diffusion & 降低计算成本 & 先压缩到潜在空间再做扩散 \\
ControlNet & 少样本条件控制 & 冻结基础模型 + 可训练编码器副本 \\
LoRA & 高效个性化 & 低秩权重更新 \\
\bottomrule
\end{tabular}
\caption{本讲核心技术总览}
\end{table}

\subsection{统一视角：Score Manipulation}

一个贯穿本讲（乃至整个扩散模型应用领域）的统一视角是 \textbf{score manipulation}——通过修改 score function（或等价的噪声预测），引导生成轨迹朝向目标方向。无论是 Classifier Guidance、CFG、还是后续课程将讨论的更多应用，其本质都是：

$$
\tilde{s}(x_t, t) = s_{\text{base}}(x_t, t) + \text{引导项}
$$

这是扩散模型区别于 GAN 和 VAE 的核心优势——生成过程的每一步都可以被外部信号灵活调控。

\subsection{工程选型建议：什么时候用哪种技术}

如果把本讲内容落到实际系统设计，可以用下面这张表快速判断技术选择：

\begin{table}[H]
\centering
\begin{tabular}{p{0.2\textwidth} p{0.24\textwidth} p{0.26\textwidth}}
\toprule
任务目标 & 优先技术 & 关键判断标准 \\
\midrule
只需要文本条件生成 & CFG / Latent Diffusion & 是否能接受一定的额外推理开销来换取更强条件一致性 \\
需要结构化控制（边缘、姿态、深度） & ControlNet & 是否已有可靠的条件信号，以及是否希望少样本快速适配 \\
需要个性化风格或主体定制 & LoRA & 是否只有少量数据、是否要求低显存和可组合部署 \\
需要保留原图结构并做编辑 & DDIM Inversion + 改条件重生成 & 是否要求对原图内容进行可控修改而非从零采样 \\
\bottomrule
\end{tabular}
\caption{Lecture 6 技术选型速查表}
\end{table}

\subsection{对后续课程的意义}

Lecture 6 的价值不只在于介绍几种具体方法，更在于建立了一个统一框架：预训练扩散模型是通用先验，而 CFG、ControlNet、LoRA、Inversion 等方法则是在不同层面把这个先验重新定向到具体任务上。后续无论是推理时引导、视频扩散还是 3D 生成，本质上都可以看成对这一框架的进一步扩展。

\subsection{拓展阅读}

\begin{itemize}
    \item \textbf{Classifier Guidance}：Dhariwal \& Nichol, ``Diffusion Models Beat GANs on Image Synthesis,'' NeurIPS 2021
    \item \textbf{Classifier-Free Guidance}：Ho \& Salimans, ``Classifier-Free Diffusion Guidance,'' NeurIPS Workshop 2022
    \item \textbf{Latent Diffusion / Stable Diffusion}：Rombach et al., ``High-Resolution Image Synthesis with Latent Diffusion Models,'' CVPR 2022
    \item \textbf{ControlNet}：Zhang et al., ``Adding Conditional Control to Text-to-Image Diffusion Models,'' ICCV 2023
    \item \textbf{LoRA}：Hu et al., ``LoRA: Low-Rank Adaptation of Large Language Models,'' ICLR 2022
    \item \textbf{DDIM}：Song et al., ``Denoising Diffusion Implicit Models,'' ICLR 2021
    \item \textbf{HuggingFace Diffusers}：\href{https://huggingface.co/docs/diffusers}{huggingface.co/docs/diffusers}——预训练扩散模型的标准库
\end{itemize}

%% --- Fin del contenido --- %%

\end{document}
